% Part: history
% Chapter: biographies
% Section: ernst-zermelo

\documentclass[../../../include/open-logic-section]{subfiles}

\begin{document}

\olfileid{his}{bio}{zer}

\olsection{এর্নস্ট জার্মেলো}

\olphoto{zermelo-ernst}{এর্নস্ট জার্মেলো}

এর্নস্ট জার্মেলোর জন্ম ১৮৭১ সালের ২৭ জুলাই
জার্মানির বার্লিনে। তাঁর পাঁচ বোন ছিলেন;
কিন্তু পরিবারে স্বাস্থ্যগত সমস্যা থাকায়
তাঁদের মধ্যে মাত্র তিনজন প্রাপ্তবয়স্ক
বয়স পর্যন্ত বেঁচে ছিলেন। জার্মেলো ছোট
থাকতেই তাঁর বাবা-মাও মারা যান; সতেরো
বছর বয়সে তিনি ও তাঁর ভাইবোনেরা
পিতামাতাহীন হয়ে পড়েন। শিল্পকলায়,
বিশেষত কবিতায়, জার্মেলোর গভীর আগ্রহ
ছিল। তিনি তীক্ষ্ণবুদ্ধি, রসিক ও
সমালোচনাপ্রবণ হিসেবে পরিচিত ছিলেন।
তাঁর সর্বাধিক খ্যাত গাণিতিক অবদানের
মধ্যে আছে ১৯০৪ সালে চয়নের
স্বতঃসিদ্ধের প্রবর্তন এবং ১৯০৮ সালে
সেটতত্ত্বের স্বতঃসিদ্ধীকরণ।

বিশ্ববিদ্যালয়ে জার্মেলোর আগ্রহ ছিল
বৈচিত্র্যময়। তিনি পদার্থবিদ্যা, গণিত
ও দর্শনের পাঠ্যক্রমে অংশ নেন।
হারমান শোয়ার্ৎসের তত্ত্বাবধানে
১৮৯৪ সালে বার্লিন বিশ্ববিদ্যালয়ে
তিনি \emph{ভ্যারিয়েশন কলন-সংক্রান্ত
অনুসন্ধান} অভিসন্দর্ভটি শেষ করেন।
১৮৯৭ সালে আরও পড়াশোনার জন্য
তিনি গ্যোটিঙেন বিশ্ববিদ্যালয়ে
যান। সেখানে ডাভিড হিলবার্টের
গণিতের ভিত্তি নিয়ে কাজ তাঁর
উপর গভীর প্রভাব ফেলে। ১৮৯৯
সালে তিনি অধ্যাপক পদের যোগ্যতা
অর্জন করেন; কিন্তু এগারো বছর
পর্যন্ত সেই পদ পাননি—সম্ভবত
তাঁর অদ্ভুত আচরণ ও “স্নায়বিক
তাড়াহুড়ো”-র কারণে।

অবশেষে ১৯১০ সালে জুরিখ
বিশ্ববিদ্যালয়ে জার্মেলো বেতনভুক্ত
অধ্যাপক পদ পান; কিন্তু যক্ষ্মার
জন্য ১৯১৬ সালে তাঁকে অবসর
নিতে হয়। সেরে ওঠার পর
১৯২১ সালে ফ্রাইবুর্গ
বিশ্ববিদ্যালয়ে তাঁকে সম্মানসূচক
অধ্যাপক পদ দেওয়া হয়। এই
সময়ে তিনি গণিতের ভিত্তি
নিয়ে কাজ করেন। থরাল্ফ
স্কোলেম ও কুর্ট গ্যোডেলের
কাজে তিনি বিরক্ত হন এবং
নিজের প্রবন্ধে প্রকাশ্যে তাঁদের
পদ্ধতির সমালোচনা করেন।
তাঁর অজনপ্রিয়তা এবং জার্মানিতে
হিটলারের ক্ষমতায় উত্থানের
বিরোধিতা—এই দুয়ের কারণে
১৯৩৫ সালে তাঁকে ফ্রাইবুর্গের
পদ থেকে সরিয়ে দেওয়া হয়।

জার্মেলোর শেষ জীবন ছিল
নিঃসঙ্গতায় চিহ্নিত। ১৯৩৫ সালে
বরখাস্ত হওয়ার পর তিনি গণিত
ছেড়ে দেন। গ্রামে চলে গিয়ে
সাদাসিধে জীবন কাটান। ১৯৪৪
সালে তিনি বিয়ে করেন; ক্রমে
দৃষ্টিশক্তি হারাতে থাকায় স্ত্রীর
উপর সম্পূর্ণ নির্ভরশীল হয়ে
পড়েন। ১৯৫১ সালের মধ্যে
তাঁর দৃষ্টিশক্তি পুরোপুরি চলে
যায়। ১৯৫৩ সালের ২১ মে
জার্মানির গুন্টারস্টালে তাঁর
মৃত্যু হয়।

\begin{reading}
জার্মেলোর পূর্ণাঙ্গ জীবনী
জানতে \citet{Ebbinghaus2015}
দেখো। তাঁর ১৯০৪ ও ১৯০৮
সালের পথিকৃৎ প্রবন্ধগুলি
মূল জার্মান ভাষায় পড়া যায়
\citep{Zermelo1904,Zermelo1908}।
পদার্থবিদ্যা-সংক্রান্ত লেখাসহ
তাঁর সংকলিত রচনাবলি ইংরেজি
অনুবাদে পাওয়া যায়
\citep{Ebbinghaus2010,Ebbinghaus2013}।
\end{reading}

\end{document}
